% Copyright 2015 Daniel Nikpayuk
\documentclass[twoside]{article}
\usepackage[letterpaper,left=2cm,right=2cm,top=2cm,bottom=2cm]{geometry}
\usepackage{amsfonts}
\usepackage{amsmath}
\usepackage{amssymb}
\usepackage{graphicx}
\usepackage{hyperref}

\newcommand{\st}{$ ^{\mbox{\scriptsize st}} $ }
\newcommand{\nd}{$ ^{\mbox{\scriptsize nd}} $ }
\newcommand{\rd}{$ ^{\mbox{\scriptsize rd}} $ }

\renewcommand{\th}{$ ^{\mbox{\scriptsize th}} $ }
%\th command exists, it gives the old english symbol used for `theta' IPA (voiceless)

\newtheorem{theorem}{Theorem}[section]
\newtheorem{lemma}[theorem]{Lemma}
\newtheorem{proposition}[theorem]{Proposition}
\newtheorem{corollary}[theorem]{Corollary}

\newenvironment{proof}[1][Proof]{\begin{trivlist}
\item[\hskip \labelsep {\bfseries #1}]}{\end{trivlist}}
\newenvironment{definition}[1][Definition]{\begin{trivlist}
\item[\hskip \labelsep {\bfseries #1:}]}{\end{trivlist}}
\newenvironment{example}[1][Example]{\begin{trivlist}
\item[\hskip \labelsep {\bfseries #1}]}{\end{trivlist}}
\newenvironment{remark}[1][Remark]{\begin{trivlist}
\item[\hskip \labelsep {\bfseries #1}]}{\end{trivlist}}
\newenvironment{note}[1][Note]{\begin{trivlist}
\item[\hskip \labelsep {\bfseries #1}]}{\end{trivlist}}

\newcommand{\qed}{\nobreak \ifvmode \relax \else
      \ifdim\lastskip<1.5em \hskip-\lastskip
      \hskip1.5em plus0em minus0.5em \fi \nobreak
      \vrule height0.75em width0.5em depth0.25em\fi}

\title{Computability Space}
\author{Daniel Nikpayuk}
\date{June 14, 2015}
\pagenumbering{alph}
\pagestyle{empty}

\begin{document}
\maketitle
\begin{figure}[h]
\centering
\includegraphics[width=1in]{../../../cc-by-nc.png}\\[0.1in]
\tiny This article is licensed under \\
\href{http://creativecommons.org/licenses/by-nc/4.0/}
{Creative Commons Attribution-NonCommercial 4.0 International.}\\[0.3in]
\end{figure}

\section*{Motivation:}

Having a background in the landscape of mathematics and some of its history, I don't think this period in time of science
and math in general is as much one for ``new discovery'' than it is for cleaning up and consolidation of existing theories.

In particular, my plan here is to take all\footnote{as problematic as a qualifier such as \emph{all} may be.} the existing
ideas of computing science and rework them into a single narrative.  What I've learned about computing science is that
computing like the \emph{humanities} tries to {\bfseries mitigate the complexity} of stories and narratives and paradigms
within their respective literatures---which is why computing for the social coder is much more connected to creative writing
than it is math.

Engineering is the originator of ``mitigating complexity'' as a \emph{way of thinking} and is certainly an inspiration to what I'm
trying to do here\footnote{I have most definitely also been inspired by ``The Structure and Interpretation of Computer Programs (SICP)}.
As a mathematician, I can honestly say that no existing branch of pure mathematics tries to mitigate complexity; each branch of math
exposes and expounds a narrative; a worldview; a perspective; a single recursive way of interacting with the world.

As the humanities have been around for a long while now, they have the intuitive framework to explicate the human experience, and so
I have borrowed my central ideas from them: Context, Semiotics, Media. I have reoriented this intuitive language as a formal structure
I call a \emph{technology space}. A way to think about it is, a technology space consists of a context space, a semiotic space, and a
media space; each being independent of each other in a formal sense; each you can think of as a branch of mathematics---each has its
own axiomatic worldview, but none by themselves mitigate complexity. It's when you take the three together, within their interplay do
you get an \emph{emergent} approximation of the complexity you're trying to model. You start with the axioms of each, and they help each
other to build each other further, allowing for a co-evolution of all three.

In particular, I aim to change the narratives of what currently exists in the landscape, recentering the philosophical idea of
``a computation''\footnote{This idea is not to be taken lightly, there are many years of research by brilliant scholars about it.}
around the coevolution of three axiomatic spaces: {\bfseries Reference, Application, Evaluation.} A reference space being the context
space, an application space being the semiotic space and an evaluation space being the media space. Such a technology space would be
called a \emph{computability space}.

Keep in mind, this idea already exists under a different package: The Lambda Calculus, which takes as its three primitives: Abstraction,
Application, and Reduction. Here, an abstraction space is the context space, an application space is again the semiotic space, and a
reduction space is the media space. The Lambda Calculus is a formal branch of applied mathematics, but focuses more on self-contained
formal proofs than it does on mitigating complexity and expressivity of computation for the general creative designer of code.

In terms of computability spaces, one actually starts out philosophically, asking what it means to ``reference''. We can play with that
idea and its connotations for a while and figure out what the minimal number of assumptions needed are (before the concept loses its
emergent meaning) and those will be our axioms for a reference space. We will play a similar game to determine the simplest natures
of our application and evaluation spaces. Finally, we can then start using these minimal spaces to help build upon each other to form
more complicated versions of each, creating new paradigms and modular forms along the way which will inevitably lead us to our mitigating
of the complexity of computability we so seek as writers of code.

\subsection*{summary}

To reiterate: A computability space is a technology space with the following bindings of undefined terms:

\begin{eqnarray}
\mbox{context space} & := & \mbox{reference} \\
\mbox{semiotic space} & := & \mbox{application} \\
\mbox{media space} & := & \mbox{evaluation}
\end{eqnarray}

To intuitively validate this against the \emph{lambda calculus}, our context is \emph{abstraction}, our semiotics is \emph{application},
while our media is \emph{reduction}.

\section*{Computability}

What does it mean to ``compute'' something? Intuitively we start with an \emph{input} and result in an \emph{output}. The part
in-between is the hidden and mysterious computation. Let's suspend our disbelief momentarily, and assume we {\bfseries do} already
know what it means to compute something, then we could immediately start looking at the broader patterns.

From this perspective, two main patterns would stand out:

\begin{enumerate}
\item	Once we were comfortable with the existence of our computations, we would start to see them as \emph{substitutions} and we
	would thus start to \emph{compose} them. It would then become a matter of becoming comfortable with the practice of taking
	the input of these compositions, substitute and resubstitute (repeatedly) until we reached a point where we had reduced our
	input to an expression where we had nothing left to substitute, in which case we would have our final output.
\item	Once we were comfortable with forming expressions of composition and evaluating their substitutions, we would start to see
	patterns within these broader initial compositions.  We would abstract out such common patterns into what would end up being
	a \emph{compressed} representation of our initial computation.  Once we became comfortable enough with these ``factorized''
	expressions, we would simply write them in the compressed form to begin with when coding new computations---the assumption
	is we'd let the {\bfseries computer} expand (or \emph{decompress}) the representation to form the tree expression of our
	composite; and then the computer would additionally perform the necessary substitutions and reductions and thus compute
	the output of interest; leaving as little tedious work to be done by the coder.
\end{enumerate}

In a way, we are looking at things backward, or we are jumping ahead, but doing so helps to build a \emph{specification} of our
simplest and initial spaces for computation. From these we can then fill in those other details of more ``realistic'' and
complicated systems of computation.

Already in this very broad approach we are running into problematics: These tree composites, how do we represent them? There's
certainly more than one way, and in the general sense no best way. Since this is the case, for the simplicity of explanation,
we will choose one approach and try to keep it consistent until we reach a competency in which we can go back and start looking
at alternatives, there benefits and costs.

The above second pattern ($ 2. $) is where things really begin to show their complexity: Let's say we stick with one representation
of composite trees: We may ourselves as coders become comfortable enough in compressing these expressions to write out ``compressed''
code from the get-go, but the truth of the matter is for many trees of composite substitutions there is more than one way to compress
(regardless of representation), and for many of those there is in no general sense a best way of doing so. It then becomes a matter
of {\bfseries modularization semantics}.

The intention then here is to explore such different strategies for compression, and when we become comfortable with this and the
existing landscape of such semantic constructs, we can then learn to apply such code based on context, optimization constraints,
and ones individual life experience.

\subsection*{References}

Getting to the core of the material then: What does it mean to ``reference'' a computation?

Any way you look at it, in its simplest form, a computation takes an \emph{input} and produces an \emph{output}. In mathematics
this is called a \emph{mapping} or a \emph{function}, but from the worldview of a computability narrative, it is better to say this
is a substitution.\footnote{a computation is an act of doing, and so is a substitution.} The ideal form of substitution is to use
a mapping with complete and instantaneously accessible information: This is to say, for whatever input we're given, when we call
our function of interest, it will immediately tell us one of two things: Is this input computable? This is to say, is this input
even associated with the mapping. If the answer is {\bfseries no}, then that in itself is the computation---this we refer to as a
\emph{null substitution}. If the answer is {\bfseries yes}, then the mapping immediately provides the output bound to the input.
A function such as this is known as a {\bfseries recursive function} in the general literature, but we ourselves will refer to it
as a {\bfseries lookup table}.

This is just a single lookup table, the reality is in practice we tend to start with several such \emph{primitive} lookup tables,
which brings up another issue. When we're given input to compute, how do we know which lookup table to call? From a \emph{meta}
perspective, one could say we will know by matter of discourse and context of the coder, which is true, but at the same time this
way of looking at it breaks the design constraint that our system of reference is as simple as possible.\footnote{such a meta-lens
assumes additional systems of reference regarding context and discourse.} As such, we will here say that the necessary information
is stored as part of the input and thus provides no redundancy of input/output bindings, in which case we can say there our reference
space is a single \emph{universal} lookup table which takes the function names as part of the input and provides the output for that
function given the input (minus the function name) applied to that function.

\subsection*{Applications}

What does it mean to ``apply'' a computation?

Our context space is a lookup table of bindings where the full input includes a recursive function name and \emph{its}
respective input.  Our semiotic space in general is a representation of this context (of bindings) which itself forms
some kind of algebra allowing for the---ideally---convenient representation of the composition of these contextual
input/output pairs.

Intuitively, if this were our specification constraint, it makes the most sense to be able to take any arbitrary collection
of \emph{finite} input and \emph{concatenate} them into another input, one which already exists within the application space.
This is to say, the basic axiom of our semiotic space is that it is \emph{self similar}.

\subsection*{Evaluations}

What does it mean to ``evaluate'' a computation in as simple a way as possible?

It is described as the simplest ideal relationship between our context and its semiotics---our media space. In particular,
any signifier of input which is itself a self similar concatenation will satisfy the very important constraint:

The output bound to this input within the lookup table matches the output given by \emph{evaluating} the component input.
This is to say, mapping each concatenated input to its respective output, and remapping the resultant output (reinterpreted
as a new input) by means of the lookup table. This process continues until no further mapping is possible, or the mapping
becomes constant.

This basic self-similar property forms our initial evaluation space.

\subsection*{An example of a finite computability space}

\section*{Expanding the context}

What would happen if we relaxed our assumption of a system of reference as function of bindings?

Regarding the how of it: As our narrative focuses on \emph{computation as an act of doing}, what if we were able to \emph{create}
new bindings within our system of reference as the computation progressed? What then would it mean to reference bindings? What if
creating new bindings reassigned input new output? Our system of reference would no longer be a function of bindings; would such
a system of computation be worth it?

From an existential point of view, assuming the full lookup table in advance, we would now be dealing with a general relation instead
of the more constrained mathematical function.  How would we keep track of the problems that arise?

\section*{\emph{fin}}

Pijariiqpunga.

\end{document}

